\subsection{A second family and a disjoint item bank}
\label{sec:family}
Two rules that we fixed before scoring decide what these runs can show. The first reads margin inflation across the 45 PolyPythias runs from a delete-one-seed jackknife with eight degrees of freedom and passes when the lower endpoint exceeds one, since nine seeds are the replicate unit and five size clusters are too few for the wild bootstrap. The second puts every bank-one item in half $A$ and every bank-two item in half $B$ on the same runs, which removes any component the two halves share only because they come from one item set. It rejects the shared-item explanation when that cross-bank interval sits above one and the jackknife interval for the within-minus-cross difference in log inflation includes zero. It supports the explanation when the difference is positive and the cross-bank interval covers one, and we report any other pattern as undecided. We simulated the first rule after fixing it, drawing seed effects from the DataDecide margin covariance and item noise from the observed reliabilities. At a true value of 1.244 it passes in 0.356 of 4,000 replicates, and at independence in 0.026, which means the rule can confirm the effect in this family but a failure can't count against it.

\outcome{\Ronecase}{1}{R1 pass}{The second family replicates the margin result. Margin inflation across the 45 runs is \pending{R1 lambda} with jackknife interval \ci{\pending{lo}}{\pending{hi}}, against 1.244 on DataDecide, and accuracy gives \pending{acc lambda} with interval \ci{\pending{lo}}{\pending{hi}}. Each PolyPythias seed is a clean rerun of one training recipe, which means the default-versus-auxiliary batch component behind our auxiliary contrast has no counterpart here. The per-size jackknives run from \pending{min} at \pending{size} to \pending{max} at \pending{size} (Appendix~\ref{app:transport}). Without BoolQ margin inflation is \pending{noBoolQ lambda} with interval \ci{\pending{lo}}{\pending{hi}}, \pending{which moves in the same direction as on DataDecide / which doesn't move the way it does on DataDecide}.}

\outcome{\Ronecase}{2}{R1 fail}{The second family doesn't replicate the margin result at the rule's threshold. Margin inflation across the 45 runs is \pending{R1 lambda} with jackknife interval \ci{\pending{lo}}{\pending{hi}}, which includes one, and accuracy gives \pending{acc lambda} with interval \ci{\pending{lo}}{\pending{hi}}. Five sizes with nine seeds each give 40 within-configuration contrasts against 250 in DataDecide, and in our simulation at a true 1.244 the middle 90\% of estimates runs from 1.026 to 1.456. A failure at this design bounds the effect in this family more than it shows the effect is absent. The per-size jackknives run from \pending{min} at \pending{size} to \pending{max} at \pending{size}, which shows where the estimate sits (Appendix~\ref{app:transport}).}

\outcome{\Rtwocase}{1}{R2 not supported}{The shared-item explanation fails its rule in this family. With bank one as half $A$ and bank two as half $B$ on the same 45 runs, margin inflation is \pending{cross lambda} with interval \ci{\pending{lo}}{\pending{hi}}, and the within-minus-cross difference in log inflation has interval \ci{\pending{lo}}{\pending{hi}}, which includes zero. A component that the halves of one item set share can't account for the inflation we measure on PolyPythias. The two banks also differ in split, in item count, and in exemplar overlap, and whatever covariance those differences carry enters the cross-bank estimate along with the run covariance.}

\outcome{\Rtwocase}{2}{R2 supported}{The shared-item explanation passes its rule in this family. With bank one as half $A$ and bank two as half $B$ on the same 45 runs, margin inflation is \pending{cross lambda} with interval \ci{\pending{lo}}{\pending{hi}}, which covers one, and the within-minus-cross difference in log inflation is positive with interval \ci{\pending{lo}}{\pending{hi}}. At least part of the within-bank inflation in PolyPythias comes from a component the two halves of one item set share, and the cross-bank value of \pending{cross lambda} is the one we'd report for its run covariance. We can't say whether DataDecide carries the same component, because its cross-bank run needs GPU time we didn't have.}

\outcome{\Rtwocase}{3}{R2 undecided}{The identification check is undecided under its rule. With bank one as half $A$ and bank two as half $B$ on the same 45 runs, margin inflation is \pending{cross lambda} with interval \ci{\pending{lo}}{\pending{hi}}, and the within-minus-cross difference in log inflation has interval \ci{\pending{lo}}{\pending{hi}}. Neither branch of the rule claims this pattern, and we report it without reading it in either direction. Bank two holds about a sixth of bank one's items, which widens both intervals.}

\subsection{Mechanism, calibration, and external checks}
Across the nine defined margin traits our estimated correlations track scoring format, averaging 0.624 within formats and 0.007 across them (Figure~\ref{fig:covariance} in Appendix~\ref{app:supplementary}). Cross-format off-diagonal covariance alone gives diagnostic inflation of 0.921 with interval \ci{0.801}{1.035} against 1.244 for the full matrix (Appendix~\ref{app:supplementary}). However, a masked covariance matrix needn't stay positive semidefinite, and each benchmark sits in exactly one format, so format and task content aren't separable here. Those two correlations and the benchmark seed variances come close to reproducing the battery. Setting every within-format pair to 0.624 and every cross-format pair to 0.007 gives 1.285 against the observed 1.244, and 1.728 against 1.786 without BoolQ. With equal variances the same correlations would give 1.803, and the gap comes from BoolQ, which holds 68\% of the trace and shares its format with no other benchmark. The fit is in sample, but re-estimating both correlations with each benchmark left out keeps every predicted deletion within 0.054 of its observed value (Appendix~\ref{app:formatfit}). A practitioner could therefore approximate inflation for a new battery from its formats and seed variances, a use we haven't tested on any other battery. Thresholding the saved margins reproduces the accuracy figure without showing how training produces large margins, since their signs alone, transformed item by item with items and runs held fixed, give 1.079 with interval \ci{0.992}{1.158} that includes one (Appendix~\ref{app:scale}). Gain adjustment doesn't reduce margin inflation, leaving it at 1.275, whereas the plan predicted more than 0.70 of the accuracy excess surviving that adjustment. The gain covariate shares item noise with the outcomes, and our shipped competence proxy is degenerate, with WinoGrande supplying more than 99.99\% of its squared standardised inputs. We therefore read none of these adjustments mechanistically (Appendix~\ref{app:proxy}).

Under independent noise our calibration simulation returns mean inflation of 1.0009 on margins and 1.0004 on accuracy. The recovery slope across effective covariance coefficients from zero to 0.50 is 0.998, and fifty item re-splits on the observed data give standard deviations of 0.002 and 0.008. These checks assess selected reference distributions and can't rule out training-schedule confounding. PolyPythias gives margin inflation of 1.406 on the three planned sizes and 1.262 with interval \ci{0.927}{1.537} on a five-size rescore, and accuracy inflation there is 1.711 with interval \ci{0.869}{2.428} and stays unbounded under test inversion (Appendix~\ref{app:transport}). Released single-configuration panels from \citet{heineman2025} give likelihood inflation of 1.705 and 1.625 on eighteen traits and accuracy inflation of 0.989 and 0.868 on eight traits, without cross-half correction, intervals, or a matched trait count. They suggest the same score-scale difference without confirming the DataDecide estimand.

\subsection{Predicting aggregate uncertainty}
\label{sec:prediction}
We find covariance gains on margins and no separation on accuracy. We compared five covariance models and an exploratory rescaled plug-in on five recipe folds, fitting the seed-covariance models off fold and giving them the held-out marginal seed standard deviations (Table~\ref{tab:prediction}). On margins the rank-one, rank-one plus gain, and full models lower mean squared error of log aggregate standard deviation by about 0.03 from the 0.0358 of independence, and their paired recipe-bootstrap intervals exclude zero. However, the rank-one interval reaches zero when nonpositive held-out variances are filled from training folds. On accuracy every such interval includes zero, and an operational comparison refitting without held-out marginals doesn't separate the full, rank-one, and shrinkage models from independence.

Under the fitted covariance the reported accuracy deviation of 0.00514 gives a sign-reversal probability of about 0.246 for a half-point gap against 0.229 under independence, and of the 1,500 observed same-size gaps, 455 fall below the corresponding 5\% risk threshold (Appendix~\ref{app:supplementary}). However, noisy gaps give no lower bound on the fraction of small true gaps. Over the 1,445 same-size recipe pairs whose margin paired difference carries a positive variance estimate, the paired-difference standard deviation has median ratio 1.024 to the independence value, with 5th to 95th percentiles of 0.39 to 1.40 (Appendix~\ref{app:paired}).

We checked the same question on real recipe comparisons. At each size below 1B we called a pair of recipes when their aggregate gap exceeded two standard errors, computed either under independence or with the estimated covariance, leaving the pair out of both estimates. We scored each call against the pair's ordering at 1B whenever that ordering itself cleared two corrected standard errors. Across the four sizes and both scales, the correction removes between 1 and 12 of the 173 to 223 calls on such pairs. That changes the share of wrong calls by at most 0.009, from rates between 0.029 and 0.070, and every recipe-bootstrap interval for that change reaches zero (Table~\ref{tab:decision}). A wrong call here also counts genuine rank changes between sizes, so these rates aren't false-positive rates in the usual sense.

\begin{table}[t]
\centering
\caption{Wrong-sign calls against the 1B ordering, as a share of calls on resolved pairs, under the independence and corrected rules. The call columns count calls on resolved pairs only, and the difference is the independence rate minus the corrected rate, with its recipe-bootstrap interval.}
\label{tab:decision}
\small\setlength{\tabcolsep}{4pt}\renewcommand{\arraystretch}{1.12}
\begin{tabular}{llrrrrr}
\toprule
& & \multicolumn{2}{c}{Calls} & \multicolumn{2}{c}{Wrong-call rate} & \\
\cmidrule(lr){3-4}\cmidrule(lr){5-6}
Score & Size & Indep. & Corrected & Indep. & Corrected & Difference\\
\midrule
Margin & 150M & 180 & 179 & 0.056 & 0.050 & \ci{0.000}{0.022}\\
Margin & 300M & 184 & 178 & 0.060 & 0.051 & \ci{-0.002}{0.049}\\
Margin & 530M & 181 & 169 & 0.061 & 0.059 & \ci{-0.016}{0.022}\\
Margin & 750M & 173 & 169 & 0.029 & 0.030 & \ci{-0.004}{0.000}\\
Accuracy & 150M & 217 & 215 & 0.069 & 0.070 & \ci{-0.003}{0.000}\\
Accuracy & 300M & 214 & 212 & 0.051 & 0.052 & \ci{-0.003}{0.000}\\
Accuracy & 530M & 223 & 218 & 0.067 & 0.064 & \ci{-0.004}{0.024}\\
Accuracy & 750M & 201 & 200 & 0.045 & 0.045 & \ci{-0.002}{0.001}\\
\bottomrule
\end{tabular}
\end{table}

